% GENERATED FILE. Edit canonical lessons, then run npm run generate:books.
% canonical-source-hash: fnv1a64:110cff520df999a0
% canonical-lessons: HI-C90-se, HI-C90-bas, HI-C90-uthna, HI-C90-kaam, HI-C90-pahle, HI-C90-baad, HI-C90-repaso-dincharya, HI-C90-sintesis-dincharya

\chapter{A Day, In The Order It Happens}
\label{ch:hi-dincharya}
\hlchaptermodality{98}

\begin{chapteropening}
\textbf{By the end of this chapter:} \emph{I can describe my daily routine in order -- what time I get up, how I travel, where I go, and what happens before and after.}

Read a four-sentence daily routine in which every word is taught, then say your own, and say it again as the other gender changing nothing but the endings.
\end{chapteropening}

\section[se]{\hi{से} — where you came from, and what you came by}

\label{lesson:HI-C90-se}



\begin{quote}
You have \textbf{\hi{पर}} for being \emph{at} a place. Try to say you came \emph{from} one, and notice you have the place and no way to leave it.
\end{quote}

\subsection*{You'll want to know: \hi{से}}
\textbf{\hi{से}} (\emph{se}) — \textquotedblleft{}from,\textquotedblright{} and also \textquotedblleft{}by.\textquotedblright{}

\textbf{\hi{घर} \hi{से}} — from home. \textbf{\hi{हाथ} \hi{से}} — by hand.

Like \textbf{\hi{पर}}, it trails its noun.

\begin{cousinweb}
English splits this into two words and Hindi does not:

\noindent
\begin{tabularx}{\linewidth}{@{}>{\raggedright\arraybackslash}X>{\raggedright\arraybackslash}X@{}}
\toprule
\textbf{English} & \textbf{Hindi} \\
\midrule
from home & \hi{घर} \hi{से} \\
by hand & \hi{हाथ} \hi{से} \\
\bottomrule
\end{tabularx}

Both are the same idea seen from two angles. \textbf{Where something starts} and \textbf{what it travels through} are, to Hindi, one relation: the thing \emph{proceeds from} the noun. A journey proceeds from your house; an action proceeds from your hand.

Once you see that, the second use stops being a second word to learn. A vehicle you travel in takes \textbf{\hi{से}} for exactly the same reason your hand does.
\end{cousinweb}

\begin{grammarlens}[title={Grammar Lens: two postpositions, one habit}]
\noindent
\begin{tabularx}{\linewidth}{@{}>{\raggedright\arraybackslash}X>{\raggedright\arraybackslash}X@{}}
\toprule
\textbf{postposition} & \textbf{what it marks} \\
\midrule
\hi{पर} & at a point \\
\hi{से} & from it, or by way of it \\
\bottomrule
\end{tabularx}

Both trail the noun, both are one syllable, and neither has an English word sitting in the same place. That is the whole class: \textbf{the small word comes second}, and knowing the noun tells you where to put it.

A question you will be asked in a speaking interview turns on this word: \emph{\hi{आप} \hi{कहाँ} \hi{से} \hi{हैं}?} — where are you from?
\end{grammarlens}

\subsection*{Your turn}
Say these aloud:

\begin{itemize}
\raggedright
  \item \textquotedblleft{}\hi{घर} \hi{से}\textquotedblright{} — from home
  \item \textquotedblleft{}\hi{हाथ} \hi{से}\textquotedblright{} — by hand
  \item \textquotedblleft{}\hi{आप} \hi{कहाँ} \hi{से} \hi{हैं}?\textquotedblright{}
  \item what the two uses have in common
\end{itemize}

\begin{tcolorbox}[breakable,colback=teal!4,colframe=teal!35!black,title={Before you move on}]
Which side of the noun does \textbf{\hi{से}} sit on? (After it.) What single idea covers \emph{from} and \emph{by}? (The thing proceeds from the noun.) Which question in a speaking interview needs it? (\emph{\hi{आप} \hi{कहाँ} \hi{से} \hi{हैं}?})
\end{tcolorbox}
\section[bas se]{\hi{बस} \hi{से} — the same three letters, twice over}

\label{lesson:HI-C90-bas}



\begin{quote}
You already read \textbf{\hi{बस}} as a reply meaning \emph{enough, that's all}. Read it again in \textbf{\hi{बस} \hi{से}} and notice it cannot possibly mean that here.
\end{quote}

\subsection*{You'll want to know: \hi{बस} \hi{से}}
\textbf{\hi{बस}} — the vehicle. \textbf{\hi{बस} \hi{से}} — by bus.

\textbf{\hi{मैं} \hi{बस} \hi{से} \hi{स्कूल} \hi{जाता} \hi{हूँ।}} — I go to school by bus.

\begin{cousinweb}
These are \textbf{two different words} that happen to be spelled alike:

\noindent
\begin{tabularx}{\linewidth}{@{}>{\raggedright\arraybackslash}X>{\raggedright\arraybackslash}X>{\raggedright\arraybackslash}X@{}}
\toprule
\textbf{word} & \textbf{where it comes from} & \textbf{what it means} \\
\midrule
\hi{बस} & Persian \emph{bas} & enough, that's all \\
\hi{बस} & English \emph{bus} & the vehicle \\
\bottomrule
\end{tabularx}

The first has been in the language for centuries. The second arrived with the vehicle, and English had shortened it from \emph{omnibus} — Latin for \textbf{\textquotedblleft{}for everybody\textquotedblright{}} — which is a fair description of a bus and a strange thing to find inside three Devanagari letters.

Nothing marks them apart on the page. \textbf{The sentence does the work}: a reply on its own is the old word, and a noun with a postposition behind it is the vehicle.
\end{cousinweb}

\begin{grammarlens}[title={Grammar Lens: reading the one from the other}]
\noindent
\begin{tabularx}{\linewidth}{@{}>{\raggedright\arraybackslash}X>{\raggedright\arraybackslash}X@{}}
\toprule
\textbf{what you see} & \textbf{which word} \\
\midrule
\hi{बस।} & enough \\
\hi{बस} \hi{से} & by bus \\
\bottomrule
\end{tabularx}

The moment \textbf{\hi{से}} follows it, it is a thing you can travel in. That is the sort of decision a reading paper tests without ever saying so, and you now make it from the shape of the sentence rather than from the word.
\end{grammarlens}

\subsection*{Your turn}
Say these aloud:

\begin{itemize}
\raggedright
  \item \textquotedblleft{}\hi{बस} \hi{से}\textquotedblright{} — by bus
  \item \textquotedblleft{}\hi{मैं} \hi{बस} \hi{से} \hi{स्कूल} \hi{जाता} \hi{हूँ}\textquotedblright{}
  \item the other word spelled the same, and what it means
  \item what tells the two apart in a sentence
\end{itemize}

\begin{tcolorbox}[breakable,colback=teal!4,colframe=teal!35!black,title={Before you move on}]
What did English shorten \emph{bus} from? (\emph{Omnibus}, Latin for \emph{for everybody}.) How do you know which \textbf{\hi{बस}} you are reading? (From what stands around it — a postposition behind it means the vehicle.)
\end{tcolorbox}
\section[uṭhnā]{\hi{उठना} — where a day starts}

\label{lesson:HI-C90-uthna}



\begin{quote}
You can say \textbf{\hi{सुबह}} and you can build a habit with \textbf{-\hi{ता}}. Try to say what you do first every morning, and find the verb missing.
\end{quote}

\subsection*{You'll want to know: \hi{उठना}}
\textbf{\hi{उठना}} (\emph{uṭhnā}) — to get up, to rise.

\textbf{\hi{मैं} \hi{सुबह} \hi{उठता} \hi{हूँ।}} — I get up in the morning. \textbf{\hi{मैं} \hi{सुबह} \hi{सात} \hi{बजे} \hi{उठती} \hi{हूँ।}} — I get up at seven in the morning.

\begin{cousinweb}
The stem is \textbf{\hi{उठ}-}, and it opens with a vowel where no consonant can carry the mark — so the word begins with the \textbf{full letter \hi{उ}}, which is why this is the first word in the track to start that way.

The verb is broader than English \emph{get up}. It covers a person leaving a bed, a person standing from a chair, smoke rising, and a price going up. What all of those share is \textbf{movement away from the ground}, and Hindi treats that as one action.

Its opposite, \textbf{\hi{बैठना}}, is \emph{to sit} — the same idea in the other direction.
\end{cousinweb}

\begin{grammarlens}[title={Grammar Lens: it takes the endings you already have}]
\noindent
\begin{tabularx}{\linewidth}{@{}>{\raggedright\arraybackslash}X>{\raggedright\arraybackslash}X@{}}
\toprule
\textbf{what you mean} & \textbf{the form} \\
\midrule
I get up (habit) & \hi{उठता} \hi{हूँ} \\
I am getting up & \hi{उठ} \hi{रहा} \hi{हूँ} \\
I used to get up & \hi{उठता} \hi{था} \\
I will get up & \hi{उठूँगा} \\
\bottomrule
\end{tabularx}

Four tenses on one new stem, and you had all four endings before you had the verb. That is what a stem costs you now: \textbf{one word, and the whole grid comes with it.}
\end{grammarlens}

\subsection*{Your turn}
\begin{itemize}
\raggedright
  \item \emph{Say it:} \textquotedblleft{}\hi{मैं} \hi{सुबह} \hi{उठता} \hi{हूँ}\textquotedblright{}
  \item \emph{Say it:} the same as a woman would say it
  \item \emph{Say it:} \textquotedblleft{}\hi{मैं} \hi{सुबह} \hi{सात} \hi{बजे} \hi{उठती} \hi{हूँ}\textquotedblright{}
  \item \emph{Look:} at \hi{उठना} and name the letter it opens with
\end{itemize}

\begin{tcolorbox}[breakable,colback=teal!4,colframe=teal!35!black,title={Before you move on}]
What do a person, smoke and a price have in common in this verb? (They all move away from the ground.) Why does the word open with a full letter? (There is no consonant to carry the mark.) What is its opposite? (\textbf{\hi{बैठना}}, to sit.)
\end{tcolorbox}
\section[kām]{\hi{काम} — the noun that was there all along}

\label{lesson:HI-C90-kaam}



\begin{quote}
You learned \textbf{\hi{काम} \hi{करना}} long ago as one expression meaning \emph{to work}. Say it, and then ask yourself what the first half of it means on its own.
\end{quote}

\subsection*{You'll want to know: \hi{काम}}
\textbf{\hi{काम}} (\emph{kām}) — work, a job, a task.

\textbf{\hi{काम} \hi{पर}} — at work. \textbf{\hi{मैं} \hi{काम} \hi{पर} \hi{जाता} \hi{हूँ।}} — I go to work.

\begin{cousinweb}
\textbf{\hi{काम} \hi{करना}} was never one word. It is \textbf{\hi{काम}} plus \textbf{\hi{करना}} — literally \emph{to do work} — and Hindi builds an enormous number of verbs this way: a noun, plus the verb \emph{to do}.

\noindent
\begin{tabularx}{\linewidth}{@{}>{\raggedright\arraybackslash}X>{\raggedright\arraybackslash}X@{}}
\toprule
\textbf{the pattern} & \textbf{what it makes} \\
\midrule
\hi{काम} + \hi{करना} & to work \\
a noun + \hi{करना} & to do that noun \\
\bottomrule
\end{tabularx}

Knowing this turns a memorised phrase into a \textbf{machine}. Any noun you meet can be handed to \textbf{\hi{करना}}, and the result is usually the verb a learner would otherwise look up separately.

The word comes from Sanskrit \textbf{\hi{कर्म}} (\emph{karma}), \textquotedblleft{}act, deed\textquotedblright{} — the same word English borrowed whole. A \emph{karma} is a thing done, and a \textbf{\hi{काम}} is a thing to do.
\end{cousinweb}

\begin{grammarlens}[title={Grammar Lens: which postposition a workplace takes}]
\noindent
\begin{tabularx}{\linewidth}{@{}>{\raggedright\arraybackslash}X>{\raggedright\arraybackslash}X@{}}
\toprule
\textbf{sentence} & \textbf{meaning} \\
\midrule
\hi{मैं} \hi{काम} \hi{पर} \hi{जाता} \hi{हूँ।} & I go to work. \\
\hi{मैं} \hi{बस} \hi{से} \hi{काम} \hi{पर} \hi{जाता} \hi{हूँ।} & I go to work by bus. \\
\bottomrule
\end{tabularx}

\textbf{\hi{पर}} for the place you end up at, \textbf{\hi{से}} for what carries you there. Both of them trail their nouns, and the verb stays at the end however many of them you stack in front of it.
\end{grammarlens}

\subsection*{Your turn}
Say these aloud:

\begin{itemize}
\raggedright
  \item \textquotedblleft{}\hi{काम}\textquotedblright{} — work
  \item \textquotedblleft{}\hi{मैं} \hi{काम} \hi{पर} \hi{जाता} \hi{हूँ}\textquotedblright{}
  \item \textquotedblleft{}\hi{मैं} \hi{बस} \hi{से} \hi{काम} \hi{पर} \hi{जाता} \hi{हूँ}\textquotedblright{}
  \item what \emph{karma} means, and how it is the same word
\end{itemize}

\begin{tcolorbox}[breakable,colback=teal!4,colframe=teal!35!black,title={Before you move on}]
What is \emph{kaam karnaa} made of? (The noun \textbf{\hi{काम}} plus \textbf{\hi{करना}}.) Which postposition marks the workplace? (\textbf{\hi{पर}}.) Which English word is the same Sanskrit root? (\emph{Karma}.)
\end{tcolorbox}
\section[pahle]{\hi{पहले} — putting one thing in front of another}

\label{lesson:HI-C90-pahle}



\begin{quote}
You can say that you get up and that you go to work. Try to say \textbf{which of the two happens first}, and notice the track has never given you a way.
\end{quote}

\subsection*{You'll want to know: \hi{पहले}}
\textbf{\hi{पहले}} (\emph{pahle}) — first, beforehand, earlier.

\textbf{\hi{पहले} \hi{मैं} \hi{उठता} \hi{हूँ।}} — First I get up.

\begin{cousinweb}
\textbf{\hi{पहला}} is \emph{first}, the ordinal, and \textbf{\hi{पहले}} is the same word in the shape an adverb takes — the \textbf{-\hi{ए}} ending you have now met on adjectives, on the past copula and on the continuous.

It does two jobs, and the difference is where it stands:

\noindent
\begin{tabularx}{\linewidth}{@{}>{\raggedright\arraybackslash}X>{\raggedright\arraybackslash}X@{}}
\toprule
\textbf{position} & \textbf{what it means} \\
\midrule
in front of the sentence & first, to begin with \\
after a time word & before that time \\
\bottomrule
\end{tabularx}

\emph{\hi{सुबह} \hi{से} \hi{पहले}} is \emph{before morning} — the time word, \textbf{\hi{से}}, then \textbf{\hi{पहले}}. That is the second use, and it is why this lesson comes after the one on \textbf{\hi{से}}.
\end{cousinweb}

\begin{grammarlens}[title={Grammar Lens: sequencing a day}]
\noindent
\begin{tabularx}{\linewidth}{@{}>{\raggedright\arraybackslash}X>{\raggedright\arraybackslash}X@{}}
\toprule
\textbf{sentence} & \textbf{meaning} \\
\midrule
\hi{पहले} \hi{मैं} \hi{उठता} \hi{हूँ।} & First I get up. \\
\hi{मैं} \hi{काम} \hi{पर} \hi{जाता} \hi{हूँ।} & I go to work. \\
\bottomrule
\end{tabularx}

Two sentences and an order between them. That is the whole of what a one-minute spoken description needs — and until this word existed, every sentence in such a description floated free of the others.
\end{grammarlens}

\subsection*{Your turn}
Say these aloud:

\begin{itemize}
\raggedright
  \item \textquotedblleft{}\hi{पहले}\textquotedblright{} — first
  \item \textquotedblleft{}\hi{पहले} \hi{मैं} \hi{उठता} \hi{हूँ}\textquotedblright{}
  \item the ordinal it comes from
  \item \textquotedblleft{}\hi{सुबह} \hi{से} \hi{पहले}\textquotedblright{}
\end{itemize}

\begin{tcolorbox}[breakable,colback=teal!4,colframe=teal!35!black,title={Before you move on}]
Which ordinal is \textbf{\hi{पहले}} built on? (\textbf{\hi{पहला}}, first.) What decides whether it means \emph{first} or \emph{before}? (Where it stands.) What must come in front of it for the second use? (A time word and \textbf{\hi{से}}.)
\end{tcolorbox}
\section[ke bād]{\hi{के} \hi{बाद} — the other end of the pair}

\label{lesson:HI-C90-baad}



\begin{quote}
You have \textbf{\hi{पहले}} for \emph{first}. Guess its partner, and then read on to see what Hindi makes you put in front of it.
\end{quote}

\subsection*{You'll want to know: \hi{के} \hi{बाद}}
\textbf{\hi{बाद}} (\emph{bād}) — after. In a sentence it almost always arrives as \textbf{\hi{के} \hi{बाद}}.

\textbf{\hi{काम} \hi{के} \hi{बाद}} — after work.

\begin{cousinweb}
\textbf{\hi{बाद}} comes from Arabic \textbf{\ar{بعد}} (\emph{baʿd}), \textquotedblleft{}after\textquotedblright{} — a plain borrowing that kept its job exactly.

What is Hindi's own is the \textbf{\hi{के}} in front of it. A two-part postposition like this one is built out of the possessive you already have: \emph{the after \textbf{of} work}. Hindi has a shelf of these, and they all share the shape:

\noindent
\begin{tabularx}{\linewidth}{@{}>{\raggedright\arraybackslash}X>{\raggedright\arraybackslash}X@{}}
\toprule
\textbf{shape} & \textbf{example} \\
\midrule
noun + \hi{के} + a word & \hi{काम} \hi{के} \hi{बाद} \\
\bottomrule
\end{tabularx}

So \textbf{\hi{के} \hi{बाद}} is not a new kind of thing. It is the possessive \textbf{\hi{का} / \hi{के} / \hi{की}} you met with \emph{uskā} and \emph{hamārā}, put to work as glue.
\end{cousinweb}

\begin{grammarlens}[title={Grammar Lens: the pair, and a day in order}]
\noindent
\begin{tabularx}{\linewidth}{@{}>{\raggedright\arraybackslash}X>{\raggedright\arraybackslash}X@{}}
\toprule
\textbf{what you mean} & \textbf{the shape} \\
\midrule
first & \hi{पहले} \\
before X & X \hi{से} \hi{पहले} \\
after X & X \hi{के} \hi{बाद} \\
\bottomrule
\end{tabularx}

Notice the asymmetry: \emph{before} takes \textbf{\hi{से}} and \emph{after} takes \textbf{\hi{के}}. There is no tidy reason to give you; it is a pair of fixed habits, and the way to keep them apart is to say each one a few times rather than to reason about it.
\end{grammarlens}

\subsection*{Your turn}
Say these aloud:

\begin{itemize}
\raggedright
  \item \textquotedblleft{}\hi{काम} \hi{के} \hi{बाद}\textquotedblright{} — after work
  \item \textquotedblleft{}\hi{काम} \hi{से} \hi{पहले}\textquotedblright{} — before work
  \item both one after the other, and hear which small word goes with which
  \item what language \textbf{\hi{बाद}} was borrowed from
\end{itemize}

\begin{tcolorbox}[breakable,colback=teal!4,colframe=teal!35!black,title={Before you move on}]
Which small word goes with \textbf{\hi{बाद}}? (\textbf{\hi{के}}.) Which goes with \textbf{\hi{पहले}}? (\textbf{\hi{से}}.) What is \textbf{\hi{के}} doing there? (It is the possessive, used as glue — \emph{the after of work}.)
\end{tcolorbox}
\section[dincaryā kī tālikā]{dincaryā kī tālikā — a day, in the order it happens}

\label{lesson:HI-C90-repaso-dincharya}



\begin{quote}
Name the verb a day starts with and the noun you travel to, before you look at the table.
\end{quote}

\begin{grammarlens}[title={Grammar Lens: the six}]
\noindent
\begin{tabularx}{\linewidth}{@{}>{\raggedright\arraybackslash}X>{\raggedright\arraybackslash}X@{}}
\toprule
\textbf{word} & \textbf{what it adds to a day} \\
\midrule
\hi{उठना} & the action it starts with \\
\hi{काम} & the place it goes to \\
\hi{बस} & what carries you there \\
\hi{से} & from, and by \\
\hi{पहले} & which of two came first \\
\hi{के} \hi{बाद} & which came second \\
\bottomrule
\end{tabularx}

Four of the six are not vocabulary in the ordinary sense. They are \textbf{joints}: things that let two sentences be put in a relation instead of standing side by side.
\end{grammarlens}

\begin{grammarlens}[title={Grammar Lens: the three small words, side by side}]
\noindent
\begin{tabularx}{\linewidth}{@{}>{\raggedright\arraybackslash}X>{\raggedright\arraybackslash}X@{}}
\toprule
\textbf{small word} & \textbf{what it marks} \\
\midrule
\hi{पर} & at a point \\
\hi{से} & from it, or by way of it \\
\hi{के} \hi{बाद} & later than it \\
\bottomrule
\end{tabularx}

All three trail their noun. That is now the whole postposition habit, and the only thing left to learn about any new one is which meaning it carries.
\end{grammarlens}

\begin{grammarlens}[title={Grammar Lens: a day, assembled}]
\noindent
\begin{tabularx}{\linewidth}{@{}>{\raggedright\arraybackslash}X>{\raggedright\arraybackslash}X@{}}
\toprule
\textbf{sentence} & \textbf{meaning} \\
\midrule
\hi{पहले} \hi{मैं} \hi{सुबह} \hi{उठता} \hi{हूँ।} & First I get up in the morning. \\
\hi{मैं} \hi{बस} \hi{से} \hi{काम} \hi{पर} \hi{जाता} \hi{हूँ।} & I go to work by bus. \\
\hi{काम} \hi{के} \hi{बाद} \hi{मैं} \hi{घर} \hi{आता} \hi{हूँ।} & After work I come home. \\
\bottomrule
\end{tabularx}

Three sentences, in order, every word of which is yours.
\end{grammarlens}

\subsection*{Your turn}
Say these aloud:

\begin{itemize}
\raggedright
  \item the three sentences in order
  \item the three small words, and what each one marks
  \item which small word goes with \textbf{\hi{पहले}} and which with \textbf{\hi{बाद}}
  \item the two words spelled \textbf{\hi{बस}}, and how a sentence tells them apart
\end{itemize}

\begin{tcolorbox}[breakable,colback=teal!4,colframe=teal!35!black,title={Before you move on}]
Which side of a noun does every postposition sit on? (After it.) Which small word does \emph{before} take, and which does \emph{after} take? (\textbf{\hi{से}} and \textbf{\hi{के}}.) What does \textbf{\hi{से}} mean besides \emph{from}? (\emph{By} — what carries you.)
\end{tcolorbox}
\section[merī dincaryā]{\hi{मेरी} \hi{दिनचर्या} — the minute you are asked to fill}

\label{lesson:HI-C90-sintesis-dincharya}



\begin{quote}
Say what time you get up, using the clock word from the future chapter. That is the first line of what follows.
\end{quote}

\begin{tcolorbox}[breakable,skin=enhanced,colback=orange!6,colframe=orange!45!black,boxrule=0.5pt,arc=1mm,left=6pt,right=6pt,top=4pt,bottom=4pt,fonttitle=\bfseries,title={Read this}]
\begin{quote}
\hi{पहले} \hi{मैं} \hi{सुबह} \hi{सात} \hi{बजे} \hi{उठती} \hi{हूँ।} \hi{मैं} \hi{बस} \hi{से} \hi{काम} \hi{पर} \hi{जाती} \hi{हूँ।} \hi{काम} \hi{के} \hi{बाद} \hi{मैं} \hi{घर} \hi{आती} \hi{हूँ।} \hi{घर} \hi{पर} \hi{मैं} \hi{हिंदी} \hi{बोलती} \hi{हूँ।}
\end{quote}

Four sentences, and \textbf{every word in them is one you have been taught}. Read it twice — once for meaning, once listening for how the joints do the work.
\end{tcolorbox}

\begin{grammarlens}[title={Grammar Lens: what is holding it together}]
\noindent
\begin{tabularx}{\linewidth}{@{}>{\raggedright\arraybackslash}X>{\raggedright\arraybackslash}X@{}}
\toprule
\textbf{joint} & \textbf{what it does} \\
\midrule
\hi{पहले} & opens the sequence \\
\hi{से} & names what carries her \\
\hi{पर} & names where she ends up \\
\hi{के} \hi{बाद} & hooks sentence three to sentence two \\
\bottomrule
\end{tabularx}

Take those four out and you have four unrelated statements. Put them back and you have a \textbf{description}, which is what the one-minute task is graded on.
\end{grammarlens}

\subsection*{How to answer}
Say your own, out loud, in four sentences:

\begin{itemize}
\raggedright
  \item what time you get up
  \item how you travel, and where to
  \item what you do after that
  \item what you do once you are home
\end{itemize}

Say it a second time and change nothing but the endings, as the other gender would. The habit endings and the joints stay exactly where they are.

\subsection*{Your turn}
Say these aloud:

\begin{itemize}
\raggedright
  \item \textquotedblleft{}\hi{पहले} \hi{मैं} \hi{सुबह} \hi{सात} \hi{बजे} \hi{उठता} \hi{हूँ}\textquotedblright{}
  \item \textquotedblleft{}\hi{मैं} \hi{बस} \hi{से} \hi{काम} \hi{पर} \hi{जाता} \hi{हूँ}\textquotedblright{}
  \item \textquotedblleft{}\hi{काम} \hi{के} \hi{बाद} \hi{मैं} \hi{घर} \hi{आता} \hi{हूँ}\textquotedblright{}
  \item all four sentences straight through, as your own day
\end{itemize}

\begin{tcolorbox}[breakable,colback=teal!4,colframe=teal!35!black,title={Before you move on}]
Which four words turn four statements into a description? (\textbf{\hi{पहले}}, \textbf{\hi{से}}, \textbf{\hi{पर}}, \textbf{\hi{के} \hi{बाद}}.) Which tense is a routine told in? (The habit, with \textbf{-\hi{ता}}.) What changes when a woman says the same day? (Only the endings.)
\end{tcolorbox}
